\begin{textsample}{POS Dim 2 – 2016\_c5 – Score -2.00 – 2016\_c5/enade\_2016\_nutri\_bdb42b9b1c.txt}  \label{ex:f2_pos_398}
CAPÍTULO 5 – DISTRIBUIÇÃO DOS CONCEITOS 5.1 - PANORAMA NACIONAL DA DISTRIBUIÇÃO DOS CONCEITOS O Conceito Enade é um indicador de qualidade que avalia o desempenho dos estudantes a partir dos resultados do Enade.
Os valores possíveis vão de 1 ( pior situação ) a 5 ( melhor situação ).
As Unidades de Observação com apenas um ou sem nenhum concluinte participante não obtêm o Conceito Enade, ficando Sem Conceito ( SC ).
A Tabela 5.1 apresenta a quantidade e distribuição de cursos de Nutrição participantes do Enade/2016, por faixa de conceito e Grande Região.
Enfatiza -se, mais uma vez, que a diferença entre os cursos tabulados neste capítulo e no capítulo 2 corresponde aos cursos sem conceito, ou seja àqueles, em princípio, sem alunos concluintes que houvessem participado da prova ou àqueles que tivessem tão somente um aluno concluinte19.
Observando -se os dados da Tabela 5.1, nota -se que, dos 326 cursos participantes, 145 ( 44,5 \% ) classificaram -se com conceito 3, o valor modal.
Esse foi também o conceito modal em quatro regiões : Nordeste ( 37,7 \% ), Sudeste ( 45,8 \% ), Sul ( 48,1 \% ) e Centro-Oeste ( 62,5 \% ).
Na região Norte, a moda foi o conceito 2, atribuído para 56,3 \% dos cursos da região.
O conceito 4 foi o segundo mais frequente em nível nacional ( 28,5 \%, correspondendo a 93 cursos ) e o conceito 2, o terceiro ( 18,7 \%, correspondendo a 61 cursos ).
Houve, ainda, 21 cursos ( 6,4 \% ) que receberam conceito 5 e quatro cursos ( 1,2 \% ) que receberam conceito 1.
Dos 326 cursos de Nutrição, dois ( 0,6 \% ) ficaram sem conceito ( SC ). 19 Estes cursos com somente um concluinte não foram considerados no capítulo 2.
A região Norte participou com 16 cursos ou 4,9 \% do total nacional.
Desses, nove cursos receberam o conceito 2, o que equivale a 56,3 \% do total regional, além de ser o conceito modal na região, como já comentado.
A um curso ( 6,3 \% ) atribuiu -se conceito 1, a quatro cursos ( 25,0 \% ), conceito 3, e a dois cursos ( 12,5 \% ), conceito 4.
Nenhum curso ficou sem conceito ( SC ) ou com conceito 5 nessa região.
A região Nordeste participou com 77 cursos ou 23,6 \% do total nacional.
Nessa região, 29 cursos ( 37,7 \% em termos regionais ) obtiveram conceito 3, o conceito modal para a região, como já comentado.
Os conceitos 1 e 2 foram atribuídos, respectivamente, a um e a 26 cursos ( 1,3 \% e 33,8 \%, respectivamente ).
Ainda nessa região, 19 cursos ( 24,7 \% ) foram avaliados com o conceito 4, e dois cursos ( 2,6 \% ), com conceito 5.
Nenhum curso ficou sem conceito ( SC ).
Dos 155 cursos participantes da região Sudeste ( 47,5 \% do total nacional ), o conceito 3 foi o conceito modal, como já anteriormente assinalado, tendo sido obtido por 71 cursos ( 45,8 \% ).
O conceito 1 foi atribuído a dois cursos ( 1,3 \% ) e o conceito 2, a 20 cursos ( 12,9 \% ).
Receberam os conceitos 4 e 5, respectivamente, 45 cursos ( 29,0 \% ) e 15 cursos ( 9,7 \% ).
Nessa região, dois cursos ( 1,3 \% ) ficaram sem conceito ( SC ).
Os 54 cursos da região Sul corresponderam a 16,6 \% do total nacional.
Como já indicado, houve predominância do conceito 3, atribuído a 48,1 \% dos cursos da região ( conceito modal ), ou seja, atribuído a 26 dos 54 cursos participantes na região Sul.
O conceito 2 foi atribuído a quatro cursos ( 7,4 \% ), o conceito 4, a 21 cursos ( 38,9 \% ) e o conceito 5, a três cursos ( 5,6 \% ).
Nenhum curso recebeu conceito 1 ou ficou sem conceito ( SC ) na região Sul.
Já dos 24 cursos participantes na região Centro-Oeste ( 7,4 \% do total nacional ), 15 receberam conceito 3, o conceito modal ( 62,5 \% ), como já fora destacado.
A dois cursos ( 8,3 \% ) foi atribuído conceito 2, e a outros seis ( 25,0 \% ), conceito 4.
O conceito 5 foi atribuído a um curso ( 4,2 \% ).
Nessa região, nenhum curso ficou sem conceito ( SC ) ou com conceito 1.
O Gráfico 5.1 apresenta a distribuição do Conceito Enade segundo Grande Região : poligonais mais à esquerda representam Grandes Regiões com pior distribuição, e poligonais mais à direita, Grandes Regiões com melhores conceitos.
A região Norte ( linha azul ) apresenta os cursos com conceitos mais baixos, ao passo que as regiões Sudeste ( linha vermelha ) e Sul ( linha roxa ) apresentam as distribuições com valores maiores.
Cursos em IES situadas nas regiões Nordeste ( linha verde ) e Centro-Oeste ( linha preta ) apresentam uma situação intermediária. 5.2 - CONCEITOS POR CATEGORIA ADMINISTRATIVA E GRANDE REGIÃO A Tabela 5.2 apresenta a distribuição dos cursos participantes do Enade/2016 de Nutrição, por Categoria Administrativa, de acordo com os conceitos por eles alcançados, segundo as Grandes Regiões brasileiras.
Dos 326 cursos participantes, 71 ( 21,8 \% ) eram ministrados em Instituições Públicas e 255 ( 78,2 \% ), em Instituições Privadas.
De acordo com as informações da Tabela 5.2, em termos nacionais, entre Instituições Públicas e Privadas, observa -se que, dos 21 cursos avaliados com conceito 5, 11 eram oferecidos em IES Públicas e dez, em IES Privadas.
Dos 71 cursos participantes de IES Públicas, o conceito 4 foi o valor modal, atribuído a 41 cursos ( 57,7 \% ).
Entre os demais cursos de Instituições Públicas participantes, dois cursos ( 2,8 \% ) foram avaliados com conceito 1, cinco cursos ( 7,0 \% da categoria ) obtiveram conceito 2, 12 cursos ( 16,9 \% ) receberam conceito 3 e, como já comentado, 11 cursos ( 15,5 \% ) receberam conceito 5.
Nesta categoria, nenhum curso ficou sem conceito ( SC ). conceito 2, 12 cursos ( 16,9 \% ) receberam conceito 3 e, como já comentado, 11 cursos ( 15,5 \% ) receberam conceito 5.
Nesta categoria, nenhum curso ficou sem conceito ( SC ).
Na rede privada, o conceito modal foi 3, atribuído a 133 cursos, o correspondente a 52,2 \% dos 255 cursos da categoria.
Entre os demais cursos participantes, dois ( 0,8 \% ) receberam conceito 1, e 56 ( 22,0 \% ), conceito 2.
O conceito 4 foi atribuído a 52 cursos ( 20,4 \% ) e o conceito 5, a dez cursos ( 3,9 \% ).
Nesta Categoria Administrativa, dois cursos ( 0,8 \% ) ficaram sem conceito ( SC ).
Na análise por região, observa -se que, na região Norte, as Instituições Privadas participaram com 12 cursos ( 75,0 \% do total regional ), a menor das proporções desta categoria entre as regiões.
Dos quais nenhum ficou sem conceito ( SC ) ou recebeu conceito 4 ou 5.
O conceito modal para as Instituições Privadas na região foi 2, atribuído a oito cursos, correspondendo a 66,7 \% dessa categoria na região.
Três cursos foram avaliados com conceito 3 e um curso, com conceito 1.
As Instituições Públicas participaram com quatro cursos na região Norte ( 25,5 \% do total regional ), a maior das proporções desta categoria nas regiões.
Dos quais dois obtiveram conceito 4, o valor modal.
Os conceitos 2 e 3 foram atribuídos a um curso cada.
Nenhum curso ficou sem conceito ( SC ) ou com conceito 1 ou 5 nessa combinação de categoria e região.
Na região Nordeste, a rede privada concentrou 58 dos 77 cursos participantes, o equivalente a 75,3 \% do total da região.
Dentre os cursos oferecidos por IES Privadas no Nordeste, 28 cursos foram avaliados com conceito 3, o conceito modal.
Um curso obteve conceito 1, 26 cursos receberam conceito 2 e três cursos alcançaram o conceito 4.
Nessa combinação de categoria e região, nenhum curso ficou sem conceito ( SC ) ou com conceito 5.
As Instituições Públicas da região Nordeste participaram com 19 cursos ( 24,7 \% do total da região ).
Desses, 16 obtiveram conceito 4, o valor modal.
Os demais cursos ficaram com conceito 3 ( um curso ) e 5 ( dois cursos ).
Nenhum curso ficou sem conceito ( SC ), e nenhum curso de IES Pública recebeu conceito 1 ou 2 nessa região.
Na região Sudeste, a proporção de cursos da rede privada foi 79,4 \%, correspondendo a 123 dos 155 cursos participantes.
Nessa categoria e região, o conceito modal foi 3, atribuído a 64 cursos.
Dois cursos ficaram sem conceito ( SC ).
Os demais foram avaliados com conceito 2 ( 17 cursos ), conceito 4 ( 30 cursos ) e conceito 5 ( dez cursos ).
Nenhum curso recebeu conceito 1.
Entre os 32 cursos oferecidos por Instituições Públicas na região Sudeste, o conceito modal foi 4, atribuído a 15 cursos.
Os demais cursos receberam os conceitos 1 ( dois cursos ), 2 ( três cursos ), 3 ( sete cursos ) e 5 ( cinco cursos ).
Na região Sudeste, nesta categoria, nenhum curso ficou sem conceito ( SC ).
As Instituições Privadas concentraram 44 dos 54 cursos participantes da região Sul, 81,5 \% do total regional, mais elevada do que a registrada nas demais regiões brasileiras.
Desses, 23 obtiveram conceito 3, o conceito modal.
Nesta combinação de Categoria Administrativa e Grande Região, os conceitos 2 e 4 foram atribuídos a quatro e a 17 cursos, respectivamente.
Nenhum curso ficou sem conceito ( SC ) ou recebeu conceito 1 ou 5.
As Instituições Públicas na região Sul participaram com dez cursos, 18,5 \% dos cursos da região.
O conceito modal foi 4, atribuído a quatro cursos.
Aos demais cursos foram atribuídos os conceitos 3 e 5 ( três cursos, cada ).
Nenhum curso da região Sul oferecido por IES Pública ficou sem conceito ( SC ) ou recebeu os conceitos 1 e 2.
Na região Centro-Oeste, 18 dos 24 cursos participantes eram de Instituições Privadas ( 75,0 \% em termos regionais ).
Desses, 15 receberam conceito 3, o conceito avaliados com conceito 3 e um curso, com conceito 1.
As Instituições Públicas participaram com quatro cursos na região Norte ( 25,5 \% do total regional ), a maior das proporções desta categoria nas regiões.
Dos quais dois obtiveram conceito 4, o valor modal.
Os conceitos 2 e 3 foram atribuídos a um curso cada.
Nenhum curso ficou sem conceito ( SC ) ou com conceito 1 ou 5 nessa combinação de categoria e região.
Na região Nordeste, a rede privada concentrou 58 dos 77 cursos participantes, o equivalente a 75,3 \% do total da região.
Dentre os cursos oferecidos por IES Privadas no Nordeste, 28 cursos foram avaliados com conceito 3, o conceito modal.
Um curso obteve conceito 1, 26 cursos receberam conceito 2 e três cursos alcançaram o conceito 4.
Nessa combinação de categoria e região, nenhum curso ficou sem conceito ( SC ) ou com conceito 5.
As Instituições Públicas da região Nordeste participaram com 19 cursos ( 24,7 \% do total da região ).
Desses, 16 obtiveram conceito 4, o valor modal.
Os demais cursos ficaram com conceito 3 ( um curso ) e 5 ( dois cursos ).
Nenhum curso ficou sem conceito ( SC ), e nenhum curso de IES Pública recebeu conceito 1 ou 2 nessa região.
Na região Sudeste, a proporção de cursos da rede privada foi 79,4 \%, correspondendo a 123 dos 155 cursos participantes.
Nessa categoria e região, o conceito modal foi 3, atribuído a 64 cursos.
Dois cursos ficaram sem conceito ( SC ).
Os demais foram avaliados com conceito 2 ( 17 cursos ), conceito 4 ( 30 cursos ) e conceito 5 ( dez cursos ).
Nenhum curso recebeu conceito 1.
Entre os 32 cursos oferecidos por Instituições Públicas na região Sudeste, o conceito modal foi 4, atribuído a 15 cursos.
Os demais cursos receberam os conceitos 1 ( dois cursos ), 2 ( três cursos ), 3 ( sete cursos ) e 5 ( cinco cursos ).
Na região Sudeste, nesta categoria, nenhum curso ficou sem conceito ( SC ).
As Instituições Privadas concentraram 44 dos 54 cursos participantes da região Sul, 81,5 \% do total regional, mais elevada do que a registrada nas demais regiões brasileiras.
Desses, 23 obtiveram conceito 3, o conceito modal.
Nesta combinação de Categoria Administrativa e Grande Região, os conceitos 2 e 4 foram atribuídos a quatro e a 17 cursos, respectivamente.
Nenhum curso ficou sem conceito ( SC ) ou recebeu conceito 1 ou 5.
As Instituições Públicas na região Sul participaram com dez cursos, 18,5 \% dos cursos da região.
O conceito modal foi 4, atribuído a quatro cursos.
Aos demais cursos foram atribuídos os conceitos 3 e 5 ( três cursos, cada ).
Nenhum curso da região Sul oferecido por IES Pública ficou sem conceito ( SC ) ou recebeu os conceitos 1 e 2.
Na região Centro-Oeste, 18 dos 24 cursos participantes eram de Instituições Privadas ( 75,0 \% em termos regionais ).
Desses, 15 receberam conceito 3, o conceito Na região Centro-Oeste, 18 dos 24 cursos participantes eram de Instituições Privadas ( 75,0 \% em termos regionais ).
Desses, 15 receberam conceito 3, o conceito modal.
Entre os demais, um recebeu conceito 2 e dois receberam conceito 4.
Nenhum curso ficou sem conceito ( SC ) ou com conceito 1 ou 5.
Quanto aos seis cursos oferecidos por Instituições Públicas na região Centro-Oeste ( 25,0 \% do total regional ), nota -se quatro cursos receberam conceito 4.
Os outros dois ficaram com conceitos 2 e 5 ( um curso, cada ).
Nenhum curso ficou sem conceito ( SC ) ou com conceito 1 ou 3.
O Gráfico 5.2 apresenta a distribuição do Conceito Enade segundo a Categoria Administrativa da IES.
Os cursos em IES Públicas ( linha azul ) apresentam uma poligonal mais à direita do que os em IES Privadas ( linha verde ) e, consequentemente, uma distribuição de Conceitos Enade com valores maiores. 5.3 - CONCEITOS POR ORGANIZAÇÃO ACADÊMICA E GRANDE REGIÃO Na Tabela 5.3, encontra -se a distribuição dos conceitos atribuídos aos cursos participantes do Enade/2016 na Área de Nutrição, por Organização Acadêmica, segundo as Grandes Regiões brasileiras.
Dos 326 cursos de Nutrição participantes, 158 eram oferecidos em Universidades, 65 em Centros Universitários e 103 em Faculdades.
Esta distribuição corresponde a, respectivamente, 48,5 \%, 19,9 \% e 31,6 \% dos cursos.
De acordo com os dados apresentados, de todos os 21 cursos avaliados com conceito 5, 20 eram vinculados a Universidades.
Os cursos vinculados a esse tipo de Organização Acadêmica tiveram o conceito 4 como conceito modal, atribuído a 68 cursos ( 43,0 \% ).
Os demais cursos vinculados a Universidades receberam os conceitos 1 ( três cursos ), 2 ( 14 cursos ), 3 ( 52 cursos ) e conceito 5 ( 20 cursos, como já mencionado ).
Ficaram sem conceito ( SC ), um único curso.
Entre os cursos em Centros Universitários, o conceito modal foi 3, atribuído a 45 cursos ( 69,2 \% ).
Nenhum curso vinculados a esse tipo de Organização Acadêmica ficou sem conceito ( SC ) ou conceito 1.
Os outros cursos receberam os conceitos 2 ( cinco cursos ), 4 ( 14 cursos ) e 5 ( um curso ).
Dos 103 cursos mantidos por Faculdades, 48 ( 46,6 \% ) receberam conceito 3, o conceito modal.
O segundo conceito mais frequentemente atribuído aos cursos oferecidos por essa categoria de IES foi o conceito 2, recebido por 42 cursos.
Na sequência, 11 cursos obtiveram conceito 4 e um curso, conceito 1.
Nenhum curso ficou com conceito 5 e um único curso ficou sem conceito ( SC ).
Considerando -se separadamente as regiões brasileiras, verifica -se que, na região Norte, as Universidades concentraram seis dos 16 cursos participantes.
Nota -se que tais cursos apresentaram distribuição bimodal, uma vez que as maiores frequências se concentraram em dois conceitos, os conceitos 3 e 4, atribuídos, cada qual, a dois cursos ( 33,3 \% cada ).
Os demais cursos receberam os conceitos 1 e 2 ( um curso, cada ).
Nenhum curso ficou sem conceito ( SC ) ou recebeu conceito 5 nesta combinação de categoria e região.
Os Centros Universitários da região Norte foram representados por três cursos, aos quais foram atribuídos os conceitos 2 ( valor modal, recebido por dois cursos ) e 3 ( recebido por um curso ).
As Faculdades participaram com sete cursos na região Norte que receberam os conceitos 2 ( seis cursos, conceito modal ) e 3 ( um curso ).
Na região Nordeste, as Universidades participaram com 26 dos 77 cursos da Área de Nutrição oferecidos na região, dos quais nenhum ficou sem conceito ( SC ).
O conceito modal foi 4, atribuído a 16 cursos, enquanto o conceito 5 foi atribuído a dois cursos.
Obtiveram os conceitos 2 e 3, respectivamente, três e cinco cursos.
Nenhum curso oferecido por Universidades no Nordeste recebeu conceito 1.
Os Centros Universitários contaram com dez cursos participante na região Nordeste.
Esses obtiveram conceitos 2 e 4 ( um curso, cada ) e 3 ( oito cursos, o conceito modal ).
Nenhum curso oferecido por Centros Universitários no Nordeste ficou sem conceito ( SC ) ou recebeu conceito 1 ou 5.
As Faculdades foram representadas por 41 cursos na região Nordeste, 22 dos quais obtiveram o conceito modal 2.
O conceito 1 foi atribuído a um curso, o conceito 3, a 16 cursos e o conceito 4, a dois cursos.
Na região Nordeste, nenhum curso oferecido por Faculdades recebeu conceito 5 ou ficou sem conceito ( SC ).
Na região Sudeste, as Universidades concentraram 84 dos 155 cursos de Nutrição da região.
Entre os cursos em Universidades, nessa região, o conceito modal foi 4, atribuído a 33 cursos.
Os demais cursos receberam os conceitos 1 ( dois cursos ), 2 ( seis cursos ), 3 ( 28 cursos ) e 5 ( 14 cursos ).
Somente um curso ficou sem conceito ( SC ).
Os Centros Universitários participaram com 34 cursos na região Sudeste, dos quais 24 obtiveram o conceito modal 3.
O conceito 2 foi atribuído a dois cursos, o conceito 4, a sete cursos, e, por fim, o conceito 5, a um curso.
Nenhum curso ficou sem conceito ( SC ) nessa combinação de Organização Acadêmica e região.
Já as Faculdades foram representadas por 37 cursos na região Sudeste, um dos quais ficou sem conceito ( SC ).
Aos demais cursos foram atribuídos os conceitos 2 ( 12 cursos ), 3 ( 19 cursos, conceito modal ) e 4 ( cinco cursos ).
Nenhum curso ficou com conceito 1 ou 5.
Dos 54 cursos da região Sul, 29 eram de Universidades, para os quais o conceito modal foi 3, atribuído a 13 cursos.
Nesse tipo de organização, nenhum curso ficou sem conceito ( SC ) ou recebeu conceito 1.
Os demais cursos receberam os conceitos 2 ( dois cursos ), 4 ( 11 cursos ) e 5 ( três cursos ).
Dos 13 cursos participantes de Centros Universitários da região Sul, sete receberam conceito 3, conceito modal, e seis, conceito 4.
Quanto aos 12 cursos vinculados a Faculdades na região Sul, seis receberam o conceito modal 3.
O conceito 2 foi atribuído a dois cursos, e o conceito 4, a quatro cursos.
Nenhum curso ficou sem conceito ( SC ) ou recebeu conceito 1, tampouco obteve conceito 5.
Na região Centro-Oeste, 13 dos 24 cursos eram de Universidades.
Nesse tipo de organização, nenhum curso recebeu conceito 1 ou ficou sem conceito ( SC ).
O conceito 2 foi atribuído a dois cursos e o conceito 3 a quatro cursos.
Os demais conceitos, 4 ( conceito modal ) e 5, foram atribuídos, respectivamente, a seis e a um curso.
Os Centros Universitários da região Centro-Oeste contaram com cinco cursos, estes obtiveram conceito 3.
Todos os seis cursos oferecidos em Faculdades na região Centro-Oeste, também, receberam conceito 3.
O Gráfico 5.3 apresenta a distribuição do Conceito Enade segundo a Organização Acadêmica da IES.
Os cursos em Universidades ( linha azul ) apresentam uma poligonal mais à direita do que as demais, denotando uma melhor distribuição de Conceitos Enade.
No outro extremo, temos os cursos oferecidos em Faculdades ( linha vermelha ) com a poligonal mais à esquerda e, consequentemente, uma distribuição de Conceitos Enade com menores valores.
IES de Centros Universitários ( linha verde ) apresentam uma situação intermediária.

% matched lemmas: 
\end{textsample}
